\section*{第 \ref{ch:b3:nervous-systems} 章——神经系统的组织}

\begin{solution}{exo:b3:nervous-systems:1}
\omterm{def:b3:nervous-systems:cells}{星形胶质细胞}：缓冲钾离子和谷氨酸，为神经元提供营养，诱导
血–脑\omterm{def:b3:innate-immunity:innate}{屏障}形成，调节局部血流。\omterm{def:b3:nervous-systems:cells}{少突胶质细胞}：
为中枢神经系统中的轴突形成\omterm{def:b3:nervous-systems:cells}{髓鞘}。施旺细胞：
为外周轴突形成\omterm{def:b3:nervous-systems:cells}{髓鞘}（并引导其再生）。\omterm{def:b3:nervous-systems:cells}{小胶质细胞}：
脑内的\omterm{def:b3:innate-immunity:innate}{巨噬细胞}，修剪突触并清除碎片。
（室管膜细胞衬于脑室内表面，并产生\omterm{prop:b3:nervous-systems:barrier-budget}{脑脊液}。）
\end{solution}

\begin{solution}{exo:b3:nervous-systems:2}
该染色使少数神经元被完全染黑，其余神经元
则不可见，因此可以看清单个细胞的全貌。卡哈尔断定，
神经元是彼此独立的细胞，相互接触而不融合，
信号沿树突到轴突的方向单向传递；高尔基则认为
它们构成一个连续的网络。卡哈尔的观点是正确的：电子
显微镜于 1954 年显示了突触间隙。
\end{solution}

\begin{solution}{exo:b3:nervous-systems:3}
轻叩股四头肌，其肌梭发放冲动；Ia
类传入纤维进入\omterm{def:b3:nervous-systems:plan}{脊髓}，直接与股四头肌运动神经元形成突触，
使之发放冲动、引起肌肉收缩，同时一个\omterm{def:b3:nervous-systems:cells}{中间神经元}
抑制股后肌群运动神经元。只有一个突触，轴突较短，
以 \qty{80}{m/s} 的速度传导，约需 \qty{30}{ms}。随意踢腿要经过
大脑——感知、决策、运动规划、下行束，
多个突触，耗时 \qtyrange{200}{300}{ms}。
\end{solution}

\begin{solution}{exo:b3:nervous-systems:4}
交感神经（靶器官处释放去甲肾上腺素）：心跳加快、心肌收缩增强，
瞳孔散大，胃肠蠕动和分泌减弱，气道扩张。
副交感神经（释放乙酰胆碱）：心跳减慢，瞳孔缩小，
胃肠蠕动和分泌增强，气道收缩。
\end{solution}

\begin{solution}{exo:b3:nervous-systems:5}
$\lambda = \sqrt{R_{m}d/4R_{i}} = \sqrt{10^{4}\times 10^{-4}/400}$ 厘米
$= \qty{0.05}{cm} = \qty{0.5}{mm}$。经过 \qty{2}{mm} 后：$e^{-4} \approx
2\,\%$。要使 $\lambda$ 加倍，直径需增至四倍，
即 \qty{4}{\micro m}。
\end{solution}

\begin{solution}{exo:b3:nervous-systems:6}
感觉神经：$12\times 6 = \qty{72}{m/s}$，$1.2/72 = \qty{17}{ms}$。运动神经：
\qty{90}{m/s}，$1.1/90 = \qty{12}{ms}$。两个突触，每个 \qty{1}{ms}，约
\qty{30}{ms}。无\omterm{def:b3:nervous-systems:cells}{髓鞘}纤维：$1.1\sqrt{12} = \qty{3.8}{m/s}$ 和
$1.1\sqrt{15} = \qty{4.3}{m/s}$: $315 + 258 + 1 \approx \qty{570}{ms}$
——太慢，无法挽救这只脚。
\end{solution}

\begin{solution}{exo:b3:nervous-systems:7}
$2.4\times 10^{20}/8.6\times 10^{10} \approx 2.8\times 10^{9}$ 个 ATP，即每个神经元每秒
消耗量；按每个动作电位耗 $2\times 10^{9}$ 个计，
平均发放频率约为 \qty{1.4}{Hz}。编码必然稀疏：发放次数少，信息由
哪些神经元在何时发放来携带，而非靠处处高频发放。
\end{solution}

\begin{solution}{exo:b3:nervous-systems:8}
两个相互抑制且不会疲劳的神经元构成一个开关：先发放的
一方会在驱动持续期间一直抑制另一方，
二者之间永远不会交替。振荡则需要一个慢过程来削弱
胜者的优势——可以是其发放的适应、
抑制性突触的压抑，或失败方在抑制解除后的反弹——该过程的
时间常数决定了周期。
\end{solution}

\begin{solution}{exo:b3:nervous-systems:9}
左旋多巴借大型中性氨基酸转运体穿过\omterm{def:b3:innate-immunity:innate}{屏障}，
中性氨基酸；多巴胺为极性且带电，无法通过。患者服用
左旋多巴（同时服用外周酶的抑制剂，以免其在穿过\omterm{def:b3:innate-immunity:innate}{屏障}前
被转化），脑内自身的酶则在需要之处由它生成多巴胺。
\omterm{def:b3:adaptive-immunity:antibody}{抗体}分子量为 \qty{150}{kDa}，既不能穿越
紧密连接，也不能借助任何转运蛋白；递送需要一种
穿梭载体劫持受体以实现转胞吞，或将药物注射入
\omterm{prop:b3:nervous-systems:barrier-budget}{脑脊液}。
\end{solution}

\begin{solution}{exo:b3:nervous-systems:10}
单位长度的膜电流由电阻性项 $V/r_{m}$ 与
电容性项 $c_{m}\,\partial V/\partial t$ 组成；电荷守恒给出
$\partial I/\partial x = -(V/r_{m} + c_{m}\,\partial V/\partial t)$，并
代入 $\partial V/\partial x = -r_{i}I$，得 $\frac{1}{r_{i}}\,\partial^{2}V/
\partial x^{2} = V/r_{m} + c_{m}\,\partial V/\partial t$；两边乘以
$r_{m}$，得 $\lambda^{2}\,\partial^{2}V/\partial x^{2} = V + \tau\,\partial
V/\partial t$。令 $\partial V/\partial t = 0$，则定理中的
方程再次得到。$\lambda$ 是长度，$\tau$ 是时间，故 $\lambda/\tau$ 是
速度——即扰动传播快慢的自然尺度。
\end{solution}

\begin{solution}{exo:b3:nervous-systems:11}
$v \propto \sqrt{d}$：速度增至四倍，需直径增至
十六倍，即 \qty{8}{mm}。截面积 $\pi(4\,\text{mm})^{2} \approx \qty{50}{mm^2}$
对比 $\pi\,(10\,\unit{\micro m})^{2} = \qty{314}{\micro m^2}$：两者之比为
$1.6\times 10^{5}$。仅由快速无\omterm{def:b3:nervous-systems:cells}{髓鞘}纤维构成的神经不可能存在
——单根此类纤维就与整条神经一样粗——因此，需要许多快速通道的动物
必须有\omterm{def:b3:nervous-systems:cells}{髓鞘}，这正是脊椎动物（以及少数无脊椎动物，
它们独立地）演化出\omterm{def:b3:nervous-systems:cells}{髓鞘}的原因。
\end{solution}

\begin{solution}{exo:b3:nervous-systems:12}
$a = 0.4/30^{0.75} = 0.4/12.8 = 0.031$；对于 \qty{70000}{g}：$0.031\times
\num{70000}^{0.75} = 0.031\times 4300 \approx \qty{134}{g}$。人类的
\qty{1350}{g} 是预测值的十倍。脑质量随体重
变化，是因为体型更大的动物需要更多感觉和运动神经元，而
神经元本身也随脑的增大而增长；与认知能力相关的
是\omterm{def:b3:nervous-systems:plan}{大脑皮层}中的神经元数量，它取决于神经元
排列得有多密——灵长类比其他哺乳动物每克脑组织中的神经元更密集，
人类则最为密集。
\end{solution}

\begin{solution}{pb:b3:nervous-systems:1}
\textbf{1.} $\lambda = \sqrt{2\times 10^{4}\,d/400}$：当 $d =
\qty{1e-4}{cm}$ 时，$\sqrt{5\times 10^{-3}} = \qty{0.071}{cm} =
\qty{0.71}{mm}$；对于 \qty{4}{\micro m}，\qty{1.41}{mm}。$\tau = 2\times
10^{4}\times 10^{-6} = \qty{20}{ms}$。
\textbf{2.} $5\,e^{-0.3/0.71} = \qty{3.3}{mV}$；$5\,e^{-0.3/1.41} =
\qty{4.0}{mV}$。
\textbf{3.} 远端输入到达时已衰减，因此突触落在何处即决定其权重；
电位约在 $\tau$ 内衰减，因此两个输入只有彼此相隔在 \qty{20}{ms} 之内
才会叠加——神经元
是一个时间窗为 \qty{20}{ms} 的符合检测器。
\textbf{4.} 有\omterm{def:b3:nervous-systems:cells}{髓鞘}者为 \qty{60}{m/s}；无\omterm{def:b3:nervous-systems:cells}{髓鞘}者 $1.1\sqrt{10} =
\qty{3.5}{m/s}$；若要使无\omterm{def:b3:nervous-systems:cells}{髓鞘}者达到 \qty{60}{m/s}，则 $d = (60/1.1)^{2}
\approx \qty{3000}{\micro m}$：即三毫米。
\textbf{5.} $\pi(5)^{2} = \qty{79}{\micro m^2}$ 与 $\pi(1500)^{2} =
7\times 10^{6}\,\unit{\micro m^2}$：在一根大轴突的截面内，
约可容纳 \num{90000} 根有髓轴突。
\textbf{6.} 以 \qty{60}{m/s} 计，\qty{2}{cm} 需 \qty{0.3}{ms}；以
\qty{3.5}{m/s} 计则需 \qty{5.7}{ms}，额外延迟约 \qty{5}{ms}。更糟的是，
裸露的膜电容大而通道少，因此来自最后一个完好郎飞结的
电流可能无法使其达到阈值：
传导阻滞。
\textbf{7.} 感觉神经：$1.2/60 = \qty{20}{ms}$；两个突触 \qty{1}{ms}；
运动神经：$1.1/90 = \qty{12}{ms}$：约 \qty{33}{ms}（另加
神经肌肉接头处的一毫秒）。
\textbf{8.} 上行至脑：$0.6/60 = \qty{10}{ms}$，加上两个突触，
信号在 \qty{20}{ms} 时进入\omterm{def:b3:nervous-systems:plan}{脊髓}，\qty{11}{ms} 之后
到达\omterm{def:b3:nervous-systems:plan}{大脑皮层}，约与脚开始移动同时；
而有意识的疼痛体验还需再多几百毫秒的加工，
因此脚在感到疼痛之前已经缩回。
\textbf{9.} $1.1\sqrt{1} = \qty{1.1}{m/s}$：$1.8/1.1 \approx
\qty{1.6}{s}$。最先出现的尖锐疼痛经细的有髓纤维传入，
在数十毫秒之内到达；钝的、烧灼样的第二次疼痛则由
无\omterm{def:b3:nervous-systems:cells}{髓鞘}的 C 纤维在一秒或更久之后传入。
\textbf{10.} 感觉神经 $3.2/60 = \qty{53}{ms}$，运动神经 $3.1/90 =
\qty{34}{ms}$，突触 \qty{1}{ms}：约 \qty{90}{ms}。身体较长的动物
反射较慢；它们靠更粗的轴突，并将更多的工作
交给局部回路来弥补。
\textbf{11.} 单程各 $0.8/80 = \qty{10}{ms}$，共 \qty{20}{ms}；突触
\qty{0.5}{ms}，接头 \qty{1}{ms}，收缩出力 \qty{5}{ms}，共 \qty{26.5}{ms}，
其余部分由肌梭自身的\omterm{def:b3:bacteriology:hgt}{转导}以及传入肌肉的传导
补足，使总数达到 \qty{30}{ms}。
\textbf{12.} 无\omterm{def:b3:nervous-systems:cells}{髓鞘}轴突的\omterm{thm:b3:nervous-systems:cable}{传导速度}仅为每秒一至两米，
因此在成体中耗时 \qty{30}{ms} 的回路需要数百毫秒，且时序欠精确；
\omterm{def:b3:nervous-systems:cells}{髓鞘}化可使速度提高十至五十倍，并使时序精确，
而行走、抓握和言语正需要这样的精确时序。
\textbf{13.} $20/50\,000 = 4\times 10^{-4}$ 摩尔/秒 $= 2.4\times 10^{20}$
个 ATP/秒；每个神经元每秒 $2.8\times 10^{9}$ 个 ATP。
\textbf{14.} $4\times 10^{8} + 7000\times 2\times 10^{4} = 4\times 10^{8}
+ 1.4\times 10^{8} = 5.4\times 10^{8}$ 个 ATP。
\textbf{15.} 若预算全部用于动作电位：$2.8\times 10^{9}/5.4\times 10^{8} \approx
\qty{5}{Hz}$；若减半则约 \qty{2.6}{Hz}。
\textbf{16.} 在 \qty{1}{Hz} 时，$5.4\times 10^{8}/2.8\times 10^{9} \approx
\qty{20}{\%}$ 的预算：编码稀疏——大多数神经元
大部分时间保持沉默，信息则由少数放电者及其放电时刻承载。
\textbf{17.} $8.6\times 10^{10}\times 50\times 5.4\times 10^{8} = 2.3
\times 10^{21}$ ATP/秒 $= 3.9\times 10^{-3}$ 摩尔/秒 $\approx \qty{190}{W}$
——约为整个静息机体的两倍。
\textbf{18.} $5\times 16 = \qty{80}{kJ/h} = \qty{22}{W}$：刚好足够，
没有余量。供应一旦中断，泵会在数秒内停止，
浓度梯度随之消散，意识在约十秒内丧失。
\textbf{19.} $a = 0.4/12.8 = 0.031$；\qty{70}{kg}：$0.031\times 4300 =
\qty{134}{g}$；\qty{5000}{kg}：$0.031\times 1.06\times 10^{5} \approx
\qty{3300}{g}$。
\textbf{20.} 人类 $1350/134 \approx 10$；大象 $4800/3300 \approx
1.5$。
\textbf{21.} 其中约 \qty{98}{\%} 位于\omterm{def:b3:nervous-systems:plan}{小脑}，负责协调
象鼻的四万块肌肉。与认知相关的是皮层神经元的数目，而非总数；
\omterm{def:b3:nervous-systems:plan}{小脑}对应的则是驱动
庞大身体运动的任务。
\textbf{22.} 随着脑增大，轴突和树突也随之延长，因此每个神经元
占据更大的体积，因而神经元数目的增长
慢于质量。在灵长类中，随着脑增大，神经元大小几乎保持不变，
因此神经元数目几乎与质量成正比增长，同样重量的灵长类脑
所含神经元数目，
是啮齿类同重量脑的数倍。
\textbf{23.} 一条腕能够依靠局部回路协调自身的吸盘和关节，伸展、
抓握、把食物沿自身传递，并避免抓住自己的皮肤；它不能看见、
选择目标，也不能学会辨别。
检验：切断腕与脑之间的神经索，再用食物触碰该腕——它仍会抓住食物，
并将其传向本应是口的位置；被截下的腕
在一小时内也会做出同样的反应。
\textbf{24.} 每个神经元约有 $23$ 个突触，相比 $7000$ 个
少了三百倍。完整的连接图谱说明哪些神经元能够影响哪些神经元，
却不能说明每个连接的强度或正负，也不能说明哪些神经调质
会改变它们，或说明动力学过程——因此即便对这种线虫而言，连接
图也只能部分预测其行为。
\textbf{25.} $\lambda = \qty{0.71}{mm}$，对于 \qty{1}{\micro m} 的
树突；缩回潜伏期约 \qty{33}{ms}；平均放电频率
约 \qty{2.6}{Hz}，其中一半预算用于动作电位。
\end{solution}
